An increase in the price of a service changes both what residents pay and what
local producers receive. These effects depend on different links: initial purchases
determine price exposure, while work and ownership determine income exposure.
We first value these changes without imposing a particular form on preferences
or technology. The aim is a local welfare expression, not yet a way to identify
the responses from observed spending.

Consider a competitive firm with smooth production $Q=F(L_0,x)$. Workplace labor
$L_0$ is fixed over the period, but the firm chooses purchased inputs $x$ after
observing the service price $p$. Their price $r$ remains fixed. Payments available
to the fixed workers and owners are the value of this restricted problem:
\[
G(p)=\max_x\{pF(L_0,x)-rx\}.
\]
At an interior optimum, $pF_x=r$: an additional input purchase raises revenue
by exactly its cost. Differentiating around the initial optimum therefore gives
\[
\mathrm dG=Q_0\,\mathrm dp+(p_0F_x-r)\,\mathrm dx=Q_0\,\mathrm dp.
\]
The input choice can change. It contributes no separate first-order term because
the initial choice is optimal. A price increase raises the payments remaining for
workers and owners by its value on initial output, rather than by the full change
in sales. Fixed labor and the given input price are important here; the expression
does not include a change in the cost of either factor.

Resident $n$ receives a fixed share $h_n$ of these payments and fixed outside
earnings $\bar v_n$, so income is $v_n=\bar v_n+h_nG(p)$. The resident chooses
service consumption $c$ and a numeraire good $m$ under $m+pc=v_n$, with general
smooth preferences. Initial optimality again removes a separate first-order
contribution from changing choices. In numeraire units, the local welfare change is
\[
\mathrm dEV_n=\mathrm dv_n-c_{n0}\,\mathrm dp
             =(h_nQ_0-c_{n0})\,\mathrm dp.
\]
Initial output values the income effect, and initial consumption values the
additional purchase cost. A resident gains from a small price increase when the
first exposure exceeds the second; living near the producer is not sufficient.

The budget line makes this comparison visible. Holding the initial production
plan fixed, a higher service price changes the line's slope while leaving its
point at $c=h_nQ_0$ unchanged. At initial consumption $c_{n0}$, the change in
available $m$ is exactly $(h_nQ_0-c_{n0})\,\mathrm dp$. The line moves outward
there when income exposure is larger, and inward when purchase exposure is larger.
For a finite change, input adjustment shifts the budget and consumers substitute;
under smoothness, these additional gains are locally second-order. The local
formula does not itself recover causal income or price changes. To use it with
spending data, the empirical analysis must supply those responses under further
demand and identifying restrictions.
